\clearpage
\begingroup
\section*{附录A：实验环境与复现配置}
\phantomsection
\addcontentsline{toc}{section}{附录A：实验环境与复现配置}
\titleformat{\subsection}{\zihao{4}\bfseries}{\thesubsection}{0.75em}{}
\renewcommand{\thesubsection}{A.\arabic{subsection}}
\renewcommand{\theHsubsection}{appendix.\arabic{subsection}}
\setcounter{subsection}{0}

\subsection{实验环境与工具版本}
\label{app:environment}
表\ref{tab:environment}记录本实验实际使用的主机、交叉编译与报告工具版本；AscendNPU 专用工具的探测信息保存在 \path{artifacts/mlir/}。
\begin{table}[H]
\centering\small
\caption{实验工具与运行环境}
\label{tab:environment}
\begin{tabularx}{\textwidth}{L{3.7cm}Y}
\toprule\rowcolor{tableHead}项目 & 实测配置 \\
\midrule
操作系统 & Windows 主机；WSL2 Ubuntu 22.04.5 LTS，x86\_64，Linux 6.6.87.2 \\
本机编译器 & GCC 11.4.0；Clang / LLVM 14.0.0 \\
交叉编译器 & riscv64-linux-gnu-gcc 11.4.0；GNU Binutils 2.38 \\
汇编实验目标 & RV64GC，LP64D，Linux ELF64 RISC-V \\
运行方式 & qemu-riscv64 6.2.0 用户态模拟 \\
辅助中间表示工具 & MLIR 14.0.0；Ascend 专用工具版本另见进阶章节及环境记录 \\
报告工具 & XeTeX / TeX Live 2025；latexmk 4.87 \\
依赖来源 & Ubuntu 22.04 官方软件仓库；SysY 运行库和 AscendNPU IR 发行包 \\
\bottomrule
\end{tabularx}
\end{table}
版本记录和已安装软件包清单见 \path{artifacts/environment/versions.txt}。表中版本来自该环境记录及原报告，不代表其他机器上的安装状态。

\subsection{目标架构与实验复现}
\label{app:riscv-environment}
实验目标为 RV64GC、LP64D 与 Linux ELF64 RISC-V。SysY 的 \texttt{int} 为 32 位，指针为 64 位，浮点参数按 LP64D 硬浮点调用约定传递\cite{riscvabi}。SysY 源程序由 Clang 14 交叉编译，手写 LLVM IR 统一使用 \texttt{llvm-as-14}、\texttt{opt-14} 和 \texttt{llc-14} 验证并生成目标文件，以保持 typed-pointer IR 与 LLVM 14 工具版本兼容；手写汇编由 \texttt{riscv64-linux-gnu-gcc} 汇编，三种实现均链接相同运行库。

生成的静态 RISC-V ELF 程序由 \texttt{qemu-riscv64} 用户态模拟器执行，用于核验指令、Linux ABI 及输入输出行为，不据此推断真实硬件性能\cite{qemu}。具体依赖、构建和测试命令见仓库 README。

\paragraph{实验源码与复现资料}
实验源码与可公开的复现资料保存在公开 GitHub 仓库：\url{https://github.com/Andy3325/compiler-principles-prelab-2026}。
固定实验快照的 Commit SHA 为
\href{https://github.com/Andy3325/compiler-principles-prelab-2026/tree/39c3bca7e140bba3b800ed7d31ddb0739b5adfb1}{\nolinkurl{39c3bca7e140bba3b800ed7d31ddb0739b5adfb1}}。
仓库包含 GCC/Clang 工具链实验、两个 SysY 示例及其手写 LLVM IR 和 RISC-V 汇编、AscendNPU IR 主机侧 lowering 实验、测试与构建脚本、报告 LaTeX 源码，以及精选结果和分析记录；README 提供目录索引、环境依赖和复现步骤。第三方源码保留原许可与来源，课程原件、校徽、二进制和重复日志不随仓库分发，完整生成产物可按 README 重建。

\Needspace{18\baselineskip}
\subsection{课程要求与实验成果对应关系}
\label{app:course-results}
表\ref{tab:results-summary}给出课程要求、实验成果及正文分析位置的对应关系。
\begin{table}[H]
\centering\small
\caption{课程要求、完成结果与对应分析}
\label{tab:results-summary}
\begin{tabularx}{\textwidth}{L{3.2cm}L{2.7cm}L{2.2cm}Y}
\toprule\rowcolor{tableHead}课程要求 & 完成情况 & 对应章节 & 实验结果与证据 \\
\midrule
基本要求一：完整编译过程 & 完成 & 第\ref{sec:toolchain}节 & 预处理、汇编文本、目标文件和链接结果；六种构建，60/60 次测试通过 \\
基本要求二：三种等价实现 & 完成 & 第\ref{sec:sysy}节 & 两个 SysY 示例，静态链接运行库；480/480 次 QEMU 执行通过 \\
进阶要求：中端 lowering & 完成主机侧中端分析 & 第\ref{sec:advanced}节 & 六次处理完成，18/18 项解析与结构检查通过 \\
进阶要求：设备端执行 & 未完成 & 第\ref{sec:advanced}节\ref{subsec:template-backend} & 缺少 hivmc/CANN 与昇腾设备，未生成设备目标文件 \\
\bottomrule
\end{tabularx}
\end{table}
\endgroup
